\documentclass{article}
\usepackage{amsmath, amssymb, amsthm}
\usepackage{hyperref}
\usepackage{geometry}
\geometry{margin=1in}

\title{Erd\H{o}s Problem \#657}
\date{Last updated: 2025-08-31}
\author{Source: erdosproblems.com}

\begin{document}
\maketitle

\noindent\textbf{Status:} OPEN \quad \textbf{No prize}

\noindent\textbf{Tags:} geometry, distances

\medskip
\noindent\textbf{Problem Statement:}

Is it true that if $A\subset \mathbb{R}^2$ is a set of $n$ points such that every subset of $3$ points determines $3$ distinct distances (i.e. $A$ has no isosceles triangles) then $A$ must determine at least $f(n)n$ distinct distances, for some $f(n)\to \infty$?

\bigskip
\noindent\textbf{Remarks:}

In \cite{Er73} Erd\H{o}s attributes this problem (more generally in $\mathbb{R}^k$) to himself and Davies. In \cite{Er97e} he does not mention Davis, but says this problem was investigated by himself, F\"{u}redi, Ruzsa, and Pach.

In \cite{Er73} Erd\H{o}s says it is not even known in $\mathbb{R}$ whether $f(n)\to \infty$. Sarosh Adenwalla has observed that this is equivalent to minimising the number of distinct differences in a set $A\subset \mathbb{R}$ of size $n$ without three-term arithmetic progressions. Dumitrescu \cite{Du08} proved that, in these terms,\[(\log n)^c \leq f(n) \leq 2^{O(\sqrt{\log n})}\]for some constant $c>0$.

Hunter observed in the comments that a result of Ruzsa coupled with standard tools of additive combinatorics (with details given by Alfaiz and Tang) allow recent progress on the size of subsets without three-term arithmetic progression (see \cite{BlSi23} which improves slightly on the bounds due to Kelley and Meka \cite{KeMe23}) yield\[2^{c(\log n)^{1/9}}\leq f(n)\]for some constant $c>0$.

Straus has observed that if $2^k\geq n$ then there exist $n$ points in $\mathbb{R}^k$ which contain no isosceles triangle and determine at most $n-1$ distances.

See also [135].

\bigskip
\noindent\textbf{References:}

[BlSi23] T. F. Bloom and O. Sisask, An improvement to the Kelley-Meka bounds on three-term arithmetic progressions. arXiv:2309.02353 (2023).
 
 [Du08] Dumitrescu, Adrian, On distinct distances and {$\lambda$}-free point sets. Discrete Math. (2008), 6533--6538.
 
 [Er73] Erd\H{o}s, P., Problems and results on combinatorial number theory. A survey of combinatorial theory (Proc. Internat. Sympos., Colorado State Univ., Fort Collins, Colo., 1971) (1973), 117-138.
 
 [Er97e] Erd\H{o}s, Paul, Some of my favourite unsolved problems. Math. Japon. (1997), 527-537.
 
 [KeMe23] Kelley, Z. and Meka, R., Strong Bounds for 3-Progressions. arXiv:2302.05537 (2023).

\bigskip
\noindent\small{Source: \url{https://www.erdosproblems.com/657}}
\end{document}
